\documentclass[11pt]{article}
\usepackage[letterpaper,margin=1in]{geometry}
\usepackage{amsmath,amssymb,amsthm,mathtools}
\usepackage{graphicx}
\usepackage{microtype}
\usepackage{xcolor}
\usepackage[hidelinks]{hyperref}

\newtheorem{theorem}{Theorem}
\newtheorem{lemma}[theorem]{Lemma}
\newtheorem{corollary}[theorem]{Corollary}
\newcommand{\Rcone}{\mathcal R_+}
\newcommand{\Ttree}{\mathcal T}
\newcommand{\N}{\mathbb N}

\title{Every Nontrivial Finite Product-Tree Inclusion Is Strict}
\author{Candidate solution packet for arXiv:1907.11476}
\date{9 August 2026}

\begin{document}
\maketitle

\noindent\fbox{\parbox{0.96\textwidth}{\textbf{Status.} Candidate full
solution, likely valid pending expert review.  The packet completely classifies
the inclusions in equation (7) of the source paper.}}

\section{Source question}

Vergara studies positive definite radial kernels on homogeneous trees and their
finite Cartesian products \cite{Vergara}.  Here $\Ttree_q$ is the
$(q+1)$-regular tree, and $\Rcone(X)$ is the cone of functions
$\varphi:\N\to\mathbb C$ for which
$(x,y)\mapsto\varphi(d_X(x,y))$ is positive definite on $X$.

For $q'\ge q$ and $N'\ge N$, isometric embeddings give
\[
 \Rcone(\Ttree_{q'}^{N'})\subset \Rcone(\Ttree_q^N). \tag{1}
\]
Immediately after equation (7), on page 6, the paper asks whether some of these
inclusions are equalities and says that a complete answer is not known.  The
source passage is reproduced below.

\begin{center}
\includegraphics[width=0.72\textwidth]{figures/open_problem_crop.png}
\end{center}

The source proves strictness when $N=N'=1$ and $q'=q+1$, and when
$N=1,N'=2,q'=q<\infty$.  It also observes that
$\Rcone(\Ttree_\infty^N)=\Rcone(\Ttree_\infty)$ for every $N$.

\section{Proof intuition}

There is no need to analyze the full moment representation.  Use the sparsest
nontrivial radial function: it is $1$ at distance zero, $t$ at distance one,
and $0$ thereafter.  On a graph $X$, its kernel operator is exactly
$I+tA_X$, where $A_X$ is the adjacency operator.

For a finite-degree product of trees,
\[
 \|A_{\Ttree_q^N}\|=2N\sqrt q.
\]
The graph is bipartite, so the bottom of the spectrum is the negative of this
number.  Hence the sparse kernel is positive definite precisely when
$|t|\le(2N\sqrt q)^{-1}$.  This threshold strictly decreases whenever either
$q$ or $N$ increases, providing a witness between every two distinct finite
parameter pairs.  If the larger tree has infinite degree, any nonzero $t$
already fails on a sufficiently large finite star.

\section{The complete classification}

We order $\{2,3,\ldots,\infty\}$ in the usual way.

\begin{theorem}[All inclusions in (1)]\label{thm:classification}
Let $q'\ge q$, $N'\ge N\ge1$, and let $q,q'$ belong to
$\{2,3,\ldots,\infty\}$.
\begin{enumerate}
\item If $q=q'=\infty$, then
\[
 \Rcone(\Ttree_{q'}^{N'})=\Rcone(\Ttree_q^N)
 =\Rcone(\Ttree_\infty).
\]
\item If $(q',N')=(q,N)$, equality is tautological.
\item In every remaining case, necessarily $q<\infty$, and the inclusion
\[
 \Rcone(\Ttree_{q'}^{N'})\subset \Rcone(\Ttree_q^N)
\]
is strict.
\end{enumerate}
For finite $q,q'$, an explicit separating function is
\[
 \varphi_t(0)=1,\qquad \varphi_t(1)=t,\qquad
 \varphi_t(n)=0\quad(n\ge2), \tag{2}
\]
where
\[
 \frac1{2N'\sqrt{q'}}<t\le \frac1{2N\sqrt q}. \tag{3}
\]
If $q'=\infty$, any $0<t\le(2N\sqrt q)^{-1}$ separates the cones.
\end{theorem}

We prove the adjacency estimate first, including the endpoint.

\begin{lemma}[Adjacency norm]\label{lem:norm}
For finite $q\ge2$ and $N\ge1$, the adjacency operator of the Cartesian
product satisfies
\[
 \bigl\|A_{\Ttree_q^N}\bigr\|=2N\sqrt q. \tag{4}
\]
Its spectrum is symmetric about zero.
\end{lemma}

\begin{proof}
First consider one tree and fix a root $o$.  Put
\[
 h(x)=\left(1+\frac{q-1}{q+1}d(o,x)\right)q^{-d(o,x)/2}.
\]
A direct check at the root and at every other vertex gives
$A_{\Ttree_q}h=2\sqrt q\,h$.  The weighted Schur estimate
\[
 2|u(x)u(y)|
 \le \frac{h(y)}{h(x)}|u(x)|^2
      +\frac{h(x)}{h(y)}|u(y)|^2
\]
summed over edges yields
$|\langle Au,u\rangle|\le2\sqrt q\,\|u\|_2^2$ for every finitely supported
$u$.  Thus $\|A_{\Ttree_q}\|\le2\sqrt q$.

For the reverse bound, let $u_R(x)=q^{-d(o,x)/2}$ for $d(o,x)\le R$ and
$u_R(x)=0$ otherwise.  Since the sphere of radius $n\ge1$ has
$(q+1)q^{n-1}$ vertices,
\[
 \|u_R\|_2^2=1+R\frac{q+1}{q},\qquad
 \langle Au_R,u_R\rangle=2R\frac{q+1}{\sqrt q}.
\]
Their quotient tends to $2\sqrt q$.

On the $N$-fold product,
\[
 A_{\Ttree_q^N}=\sum_{j=1}^N
 I\otimes\cdots\otimes A_{\Ttree_q}\otimes\cdots\otimes I.
\]
The triangle inequality gives the upper bound $2N\sqrt q$.  The tensor vectors
$u_R^{\otimes N}$ have Rayleigh quotients equal to $N$ times the one-tree
quotient, proving equality in (4).  Finally, multiplication by
$(-1)^{\sum_j d(o,x_j)}$ conjugates the adjacency operator to its negative,
so the spectrum is symmetric.
\end{proof}

\begin{lemma}[Exact sparse-kernel threshold]\label{lem:sparse}
For finite $q$ and $N\ge1$, the function in (2) belongs to
$\Rcone(\Ttree_q^N)$ if and only if
\[
 |t|\le\frac1{2N\sqrt q}. \tag{5}
\]
For every $N\ge1$ and every nonzero $t$, it does not belong to
$\Rcone(\Ttree_\infty^N)$.
\end{lemma}

\begin{proof}
For any graph $X$, the kernel associated with (2) is entrywise
\[
 K_t=I+tA_X. \tag{6}
\]
For a finite-degree product, positive definiteness of the kernel is equivalent
to nonnegativity of (6) on finitely supported vectors.  Lemma \ref{lem:norm}
and spectral symmetry show that this holds exactly when
$1-|t|\,2N\sqrt q\ge0$, proving (5).  Equivalently, failure above the threshold
can be seen directly by applying the parity sign to the tensor approximate
extremizers in the proof of Lemma \ref{lem:norm}.

For $\Ttree_\infty^N$, fix all but one coordinate and select a vertex together
with $M$ distinct neighbors in the remaining coordinate.  The corresponding
principal matrix is $I+tA_{K_{1,M}}$, whose eigenvalues include
$1-|t|\sqrt M$.  This is negative for sufficiently large $M$ whenever
$t\ne0$.
\end{proof}

\begin{proof}[Proof of Theorem \ref{thm:classification}]
The infinite-degree equality is Proposition 1.5 and the paragraph following it
in the source \cite{Vergara}; identical parameters give the second case.

Suppose now that $q<\infty$ and the parameter pairs are distinct.  If $q'$ is
finite, monotonicity and at least one strict parameter inequality give
\[
 2N'\sqrt{q'}>2N\sqrt q.
\]
Choose $t$ as in (3).  Lemma \ref{lem:sparse} puts $\varphi_t$ in the cone on
the right of (1) but not in the cone on the left, so the inclusion is strict.
If $q'=\infty$, choose any positive $t$ satisfying the finite threshold for
$(q,N)$; the second part of Lemma \ref{lem:sparse} excludes it from the larger
product.  These cases exhaust (1).
\end{proof}

\section{A slightly more general separation}

The same proof gives
\[
 \left\|A_{\Ttree_{q_1}\times\cdots\times\Ttree_{q_N}}\right\|
 =2\sum_{j=1}^N\sqrt{q_j} \tag{7}
\]
when all $q_j$ are finite.  Thus the sparse witness separates any coordinatewise
proper inclusion between anisotropic finite products as well.  This is not
needed for Theorem \ref{thm:classification}, but shows that the mechanism is
intrinsic to products rather than to equal factor degrees.

\section{Verification, limitations, and novelty}

The proof uses only finite-support quadratic forms.  In particular, it does not
depend on the product moment representation or on convergence of infinite
matrices.  The endpoint in (5) is included by the explicit upper bound, while
failure beyond it is witnessed by finitely supported approximate extremizers.
The optional script \texttt{code/verify\_sparse\_witness.py} checks the displayed
Rayleigh formulas and sample separating parameters; it is not part of the
proof.

A bounded literature search on 9 August 2026 used the run's cheap indexes, the
exact title and arXiv id, the author name, the displayed inclusion, and the
phrases ``positive definite radial kernels,'' ``products of homogeneous
trees,'' ``strict inclusion,'' and ``adjacency operator.''  It found the source
and related radial Schur-multiplier background, but no later result giving this
classification or the witness (2).  Novelty confidence is moderate pending
specialist review.

\section{Human-review recommendation}

The packet is ready for expert review as a candidate full answer to equation
(7).  The decisive checks are the weighted Schur estimate in Lemma
\ref{lem:norm}, the tensor-product lower bound, and the equivalence (6) at the
endpoint.  The literature status should also be checked by a specialist in
radial multipliers on trees.

\begin{thebibliography}{9}
\bibitem{Vergara}
I. Vergara,
\emph{Positive definite radial kernels on homogeneous trees and products},
arXiv:1907.11476v2 (2019); J. Operator Theory \textbf{84} (2020), 435--460.
\end{thebibliography}

\end{document}
